\documentclass[10pt,a4paper]{amsart}
\usepackage[margin=19mm]{geometry}
\usepackage{amsmath,amssymb,mathtools}
\usepackage[scr=rsfs,cal=euler]{mathalfa}
\usepackage{xfrac}
\usepackage[unicode,hidelinks,bookmarks=false]{hyperref}
\newcommand{\ie}{i.e.\ }
\newcommand{\cf}{cf.\ }
\newcommand{\resp}{respectively\ }
\newcommand{\Subsection}{\subsection}
\providecommand{\nobreakdash}{\nobreak\hbox{-}}
\setlength{\emergencystretch}{2em}
\multlinegap0pt
\usepackage{bm}
\usepackage{bbm}
\usepackage{dsfont}
\usepackage{xspace}
\usepackage{cancel}
\usepackage{mathtools}
\usepackage{amsthm}
\makeatletter
\@ifundefined{theorem}{\newtheorem{theorem}{Theorem}}{}
\makeatother
\title[Labeled Chip-Firing on Undirected $k$-ary Trees]{Labeled Chip-Firing on Undirected $k$-ary Trees}
\author{Ryota Inagaki, Aaron Lin}
\date{Source: arXiv:2509.17358v2 (CC BY 4.0)}
\begin{document}
\maketitle

\noindent\emph{Source and licence.} Labeled Chip-Firing on Undirected $k$-ary Trees. Version arXiv:2509.17358v2 (CC BY 4.0), distributed under the Creative Commons Attribution 4.0 International licence, as recorded in the arXiv licence metadata for this exact version and shown on the abstract page https://arxiv.org/abs/2509.17358v2. Full rights notice: licenses/arxiv-2509.17358-CC-BY-4.0.txt.\par\smallskip

\noindent\textit{Editorial scope.} Counted unit: the Theorem whose source label is \texttt{thm:Endgame} in the article (source0.tex lines 241--248, source label \texttt{thm:Endgame}), delivered together with the article's own complete original proof of it (source0.tex lines 250--278, one \texttt{proof} environment of 29 lines). No mathematics is written, repaired or re-derived: statement and proof are the article's own text line for line. Required context retained uncounted: none. Every definition, display and label of those lines is retained unchanged. Omitted from this package and not redistributed: 1--240, 249--249, 279--833. Nothing inside a retained excerpt is omitted or altered: every retained body is delivered exactly as the article's lines are written, so no omission receipt of any kind accompanies this package. Citations: the retained excerpts carry no citation command at all, so no bibliography entry of the article is transcribed and no reference is redistributed. Reference adaptation: none was needed and none was made; every reference inside the delivered unit resolves inside it, so no unresolved reference or citation is produced. Printed slips kept literal: no printed word, symbol, hypothesis or number of the counted statement or of its proof was altered. Layout and class: the article's own class is \texttt{article} with packages including amsmath, amssymb, amscd, amsthm, amsfonts, graphicx, hyperref, float, comment, inputenc; the wrapper is the lane's pinned \texttt{amsart} editorial wrapper at 10pt on A4, which restates the theorem-like environments the retained text needs and adds the article's own macro definitions that the retained text needs (none), with extra wrapper packages (bm, bbm, dsfont, xspace, cancel, mathtools). The delivered numbering is the unit's own. Assets: no retained span contains a figure environment, a graphics-inclusion command, a PDF-page inclusion, a table or a TikZ picture; no image file is copied into this package, the package declares no asset dependency and no image file is redistributed with it. Licence: version arXiv:2509.17358v2 of this article is distributed under the Creative Commons Attribution 4.0 International licence, as recorded in the arXiv licence metadata for this exact version and shown on the abstract page https://arxiv.org/abs/2509.17358v2. A full rights notice is delivered at licenses/arxiv-2509.17358-CC-BY-4.0.txt. Characters: every retained excerpt is pure ASCII, so the delivered engine, XeLaTeX with its default Latin Modern text font, reports no missing character.
\begin{theorem}\label{thm:Endgame}
Consider a vertex $i$ with $\mathrm{layer}(i) > 1$. Then for all $j$ such that $\mathrm{layer}(i)+j < \ell$, we have:

\begin{itemize}
    \item $\left( \left \lfloor \frac {i-1} k\right \rfloor, j \right) < (i,j) < \left( \left \lfloor \frac {i-1} k\right \rfloor, j+1 \right)$
    \item $(i,j)$ occurs when vertex $i$ has $k+1$ chips.
\end{itemize}
\end{theorem}

\begin{proof}
We prove the above proposition by induction on both $\mathrm{layer}(v_i)$ and $j$. We do this by proving the following steps:
\begin{enumerate}
    \item We have $(0,0) < (a, 0)$ for all $a \in \{1, 2, \dots, k\}$.
    \item We have $(i,0) < ({ki+a},0)$ for all $a \in \{1, 2, \dots, k\}$, assuming $\left ( \left \lfloor \frac {i-1} k \right \rfloor , 0 \right ) < (i,0)$.
    \item For all $a \in \{1, 2, \dots, k\}$, we have $(ki+a, 0) < (i,1)$.
    \item We have $(0,j) < (a, j)$ for all $a \in \{1, 2, \dots, k\}$, assuming inductively $(a',j-1)< (0,j)$ for all $a' \in \{1, 2, \dots, k\}$.
    \item We obtain $(i,j) < ({ki+a'}, j)$ for all $a' \in \{1,2,\dots, k\}$, assuming inductively $\left ( \left \lfloor \frac {i-1} k \right \rfloor , j \right ) < (i,j)$ and $(ki+a, j-1) < (i,j)$ for all $a \in \{1,2, \dots, k\}$
    \item We have $(i,j) < \left ({\left \lfloor \frac {i-1} k \right \rfloor}, j+1 \right)$, assuming inductively $({ki+a}, j-1) < (i,j)$ for all $a \in \{1,2, \dots, k\}$.
\end{enumerate}

When steps 1, 2, and 3 are proven, we have proven the base case. Proving steps 4, 5, and 6 proves the inductive step. 

We now prove each step of the induction.

\begin{enumerate}
    \item Suppose $(0,0)$ occurs before $(1,0)$. Then, after $(0,0)$, vertex 1 will have at least 1 chip. Further, $(1,0)$ will then increase the number of chips at $(0,0)$ by another chip, so it will have at least 2 chips at the end, a contradiction. Similarly, we know that $(0,0) < (a,0)$ for all $a \in \{2, 3, \dots, k\}$. Therefore, $(0,0)$ must occur when vertex 0 has $k+1$ chips, as it needs to end with 1 chip.

    \item Suppose $(i,0)$ occurs before $(ki+1, 0)$. Then after $(i,0)$ occurs, vertex $i$ will gain two chips from $(ki+1, 0)$ and $\left ( \left \lfloor \frac {i-1} k \right \rfloor, 0 \right )$, a contradiction. Similarly, we know that $(i,0) < (ki+a,0)$ for all $a \in \{2, 3, \dots, k\}$. Further, since vertex $i$ must end with 1 chip, $(i,0)$ must occur when vertex $i$ has $k+1$ chips.

    \item Suppose $(ki+1,0)$ occurs before $(i,1)$. Then after $(ki+1, 0)$, vertex $ki+1$ will gain two chips from $(i,1)$ and $(i,0)$, a contradiction. Therefore, $(ki+a, 0) < (i, 1)$ for all $a \in \{1, 2, \dots, k\}$.

    \item Suppose $(0,j)$ occurs before $(1,j)$. Then since $(s,j-1)< (0,j)$ for all $s \in \{1,2, \dots, k\}$, after $(0,j)$, vertex 0 will keep at least 1 chip, lose $k(j-1)$, gain at least $j$ chips from vertex 1, and gain at least $(k-1)(j-1)$ chips from vertices $2, 3, \ldots, k$. Then vertex 0 will have at least 2 chips, a contradiction. Similarly, we know that $(0,j) < (a,j)$ for all $a \in \{1, 2, \dots, k\}$. Therefore, after $(0,j)$, vertex 0 will lose $k(j-1)$ chips and gain $k(j-1)$, meaning $(0,j)$ must occur when vertex 0 has $k+1$ chips.

    \item Suppose $(i,j)$ occurs before $(ki+1, j)$. Then after $(i,j)$ vertex $i$ will lose $(k+1)(j-1)$ chips. Further, since $\left ( \left \lfloor \frac {i-1} k \right \rfloor, j \right ) < (i,j)$, vertex $i$ will receive at least $j$ chips from vertex $\left \lfloor \frac {i-1} k \right \rfloor$, $j$ chips from vertex $ki+1$, and $(k-1)(j-1)$ chips from vertices $ki+2, \ldots, ki+k$. Then vertex $i$ will have at least 2 chips at the end, a contradiction. Further, since $(i,j) < (ki+a', j)$ for all $a' \in \{1,2, \dots, k\}$, after $(i,j)$, vertex $i$ will lose $(k+1)(j-1)$ chips and gain $j+k(j-1) = (k+1)j-k$ chips, so $(i,j)$ must occur when vertex $i$ has $k+1$ chips.

    \item Suppose $(i,j)$ occurs before $\left(\left \lfloor \frac {i-1} k \right \rfloor, j+1 \right)$. Then after $(i,j)$, vertex $i$ will lose $(k+1)(j-1)$ more chips. However, since $(ki+a, j-1) < (i,j)$ for all $a \in \{1, 2, \dots, k\}$, vertex $i$ will gain at least $k(j-1)$ chips from vertices $ki+1, \ldots, ki+k$ and also gain $j+1$ chips from vertex $\left \lfloor \frac {i-1} k \right \rfloor$. Therefore, vertex $i$ receives at least $(k+1)j-1$ chips, leaving vertex $i$ with at least 2 chips at the end, a contradiction. Therefore, $(i,j) < \left(\left \lfloor \frac {i-1} k \right \rfloor, j+1 \right).$
\end{enumerate}
\end{proof}

\end{document}
